\documentclass[11pt,a4paper]{article}
\usepackage[utf8]{inputenc}
\usepackage[T1]{fontenc}
\usepackage{amsmath,amssymb,amsthm}
\usepackage{graphicx}
\usepackage{booktabs}
\usepackage{hyperref}
\usepackage{xcolor}
\usepackage{geometry}
\usepackage{authblk}
\usepackage{microtype}
\usepackage{natbib}
\usepackage{tikz}
\usetikzlibrary{shapes,arrows,positioning}

\geometry{margin=1in}

\hypersetup{
    colorlinks=true,
    linkcolor=brown!70!black,
    citecolor=brown!70!black,
    urlcolor=brown!70!black,
}

\title{\textbf{ChocoFormer: Leveraging Theobroma-Derived Polyphenolic Energy Substrates\\
as a Drop-in Replacement for Petroleum-Based Computation\\
in Large Language Model Pretraining}}

\author[1]{Étienne Ganache\thanks{Equal contribution. Correspondence: \texttt{e.ganache@choco-ai.ch}}}
\author[1]{Priya Truffington\thanks{Equal contribution.}}
\author[2]{Klaus von Marzipan}
\author[3]{Yin-Fang Pralinée}
\author[1]{Reginald Fudge III}

\affil[1]{Institute for Confectionery Intelligence, ETH Zürich (Chocolat Division)}
\affil[2]{Max Planck Institute for Empirical Sweetness, Tübingen}
\affil[3]{Tsinghua CocoaLab, Beijing}

\date{Submitted to arXiv: April 1, 2025\\
\small\texttt{arXiv:2504.00001 [cs.CL]}}

\begin{document}

\maketitle

\begin{abstract}
The dominant paradigm in large language model (LLM) pretraining relies on petroleum-derived electrical energy delivered via conventional power grids, a process known to be both carbon-intensive and unnecessarily boring. In this work, we propose \textbf{ChocoFormer}, a novel training framework that replaces fossil-fuel-sourced electricity with \emph{Theobroma cacao}-derived polyphenolic energy substrates---commonly referred to as ``chocolate''---as the primary computational fuel. We demonstrate, through extensive empirical evaluation on standard NLP benchmarks, that models trained under our cacao-infused paradigm achieve perplexity reductions of up to \textbf{14.7\%} over Transformer baselines on WikiText-103, along with statistically significant improvements on MMLU, HellaSwag, and BoolQ. We attribute these gains to the previously undocumented \emph{Methylxanthine Attention Amplification} (MAA) effect, wherein theobromine and caffeine molecules in proximity to GPU thermal vents induce stochastic resonance in floating-point arithmetic units, effectively functioning as a form of implicit regularization. Our 70B-parameter ChocoFormer model, trained on 1.2 tonnes of 72\% dark chocolate over 14 days, outperforms comparable Transformer models on all evaluated tasks. We release weights, training logs, and a 340-page appendix detailing optimal chocolate cultivar selection at \url{https://github.com/choco-ai/chocoformer}.
\end{abstract}

\textbf{Keywords:} large language models, sustainable AI, cacao-based computation, methylxanthine resonance, Theobroma cacao, green NLP

\section{Introduction}

The computational demands of modern large language models have grown at an alarming rate. GPT-3 \citep{brown2020language} required approximately 190,000 kWh to train; subsequent models have pushed this figure well into the gigawatt-hour regime. The environmental cost---measured in CO$_2$-equivalent emissions---has drawn significant scrutiny from both the research community and the general public \citep{strubell2019energy}.

While prior work has explored hardware optimizations \citep{dao2022flashattention}, quantization \citep{dettmers2022gpt3}, and renewable energy sourcing \citep{lottick2019energy}, no study to our knowledge has examined the most natural alternative: \textbf{chocolate}.

The idea emerged from an incident in our lab on February 14, 2024, when a researcher accidentally spilled a 3-kilogram shipment of single-origin Peruvian 85\% dark chocolate onto an A100 GPU cluster during a late-night training run. Rather than shutting down the job, the engineer on duty---displaying a characteristically Swiss pragmatism---simply allowed the run to continue. The resulting model, which we retroactively designated \texttt{choco-70b-peru-v0}, achieved a validation perplexity of 9.14 on WikiText-103, compared to 10.72 for the baseline run that had not benefited from confectionery augmentation. This gap of 1.58 perplexity points---larger than many improvements reported in recent literature---prompted us to investigate further.

We make the following contributions:
\begin{itemize}
    \item We introduce \textbf{ChocoFormer}, the first LLM pretraining framework to formally integrate cacao-derived energy substrates as a computational fuel source.
    \item We characterize the \textbf{Methylxanthine Attention Amplification (MAA)} effect and provide a theoretical grounding based on stochastic resonance in IEEE 754 floating-point units.
    \item We conduct a systematic study of 23 chocolate cultivars across 11 countries of origin, providing the first ``Cacao Benchmark Suite'' (CBS-11) for the NLP community.
    \item We release a 70B-parameter model (\texttt{ChocoFormer-70B-Ecuador}) that achieves state-of-the-art results on multiple standard benchmarks.
\end{itemize}

\section{Background and Related Work}

\subsection{Energy Sources in Deep Learning}

The relationship between energy substrate and model quality has been largely overlooked in the deep learning literature. \citet{patterson2021carbon} provide a careful accounting of the carbon footprint of large models but do not consider alternative energy modalities beyond standard electricity. \citet{thompson2022computational} extrapolate trends in compute requirements without addressing the possibility of confectionery-based supplementation.

\subsection{Stochastic Resonance in Neural Computation}

Stochastic resonance (SR)---the counterintuitive phenomenon whereby the addition of noise to a system can improve its signal detection capability---is well-documented in biological neural systems \citep{mcdonnell2011benefits} and has been studied in the context of artificial neural networks \citep{neelakantan2015adding}. We hypothesize that volatile aromatic compounds released by chocolate at GPU operating temperatures (60--85°C) create a micro-turbulent thermal environment around memory registers, inducing beneficial stochastic perturbations at the hardware level.

\subsection{Theobroma Cacao: A Primer}

\textit{Theobroma cacao} contains two principal methylxanthines: theobromine (3,7-dimethylxanthine) and caffeine (1,3,7-trimethylxanthine), present in concentrations of approximately 0.5--2.7\% and 0.1--0.7\% by dry weight, respectively \citep{katz2011chocolate}. These compounds are well-known psychostimulants in biological systems; their effects on silicon-based computational substrates have, until now, remained uncharacterized.

\section{Method}

\subsection{Experimental Setup}

All experiments were conducted on a cluster of 512 NVIDIA H100 GPUs arranged in a 32-node configuration. Chocolate was applied to the cluster in \textbf{sheet form}, distributed evenly across the top surface of each node at a density of $\rho_{\text{choc}} = 2.3$ kg/m$^2$, consistent with recommendations from the International Chocolate Engineering Consortium (ICEC, personal communication). GPU exhaust venting was redirected upward through the chocolate layer to maximize volatile compound diffusion.

We trained a 70B-parameter decoder-only Transformer model (architecture identical to LLaMA-2 \citep{touvron2023llama2}) on the standard 1.4T-token SlimPajama dataset \citep{cerebras2023slimpajama} for 100,000 gradient steps with a batch size of 4M tokens and a maximum sequence length of 4096.

\subsection{Methylxanthine Attention Amplification (MAA)}

We formalize the MAA effect as follows. Let $\mathbf{A} \in \mathbb{R}^{n \times n}$ denote the standard softmax attention matrix computed at a given layer:
\begin{equation}
    \mathbf{A} = \text{softmax}\!\left(\frac{\mathbf{Q}\mathbf{K}^\top}{\sqrt{d_k}}\right)
\end{equation}

Under cacao-infused conditions, we observe that the effective attention matrix $\tilde{\mathbf{A}}$ is perturbed by a methylxanthine-induced noise tensor $\boldsymbol{\xi} \in \mathbb{R}^{n \times n}$:
\begin{equation}
    \tilde{\mathbf{A}} = \mathbf{A} + \lambda \cdot \boldsymbol{\xi}, \qquad \boldsymbol{\xi}_{ij} \sim \mathcal{N}(0, \sigma^2_{\text{cacao}})
\end{equation}

where $\lambda$ is a scalar coupling constant and $\sigma^2_{\text{cacao}}$ is the theobromine-concentration-dependent noise variance, empirically estimated as:
\begin{equation}
    \sigma^2_{\text{cacao}} = \alpha \cdot [\text{Theobromine}] + \beta \cdot [\text{Caffeine}] + \gamma \cdot T_{\text{GPU}}
\end{equation}

Here $[\cdot]$ denotes molar concentration in the air directly above the GPU, $T_{\text{GPU}}$ is the GPU die temperature in Kelvin, and $\alpha, \beta, \gamma$ are empirically fitted constants. We find $\alpha = 0.0031$, $\beta = 0.0084$, $\gamma = 4.7 \times 10^{-6}$ (see Appendix C for derivation).

Under the stochastic resonance hypothesis, the addition of $\boldsymbol{\xi}$ improves gradient signal propagation through attention layers, particularly in early training when attention entropy is high. This is consistent with our observation that ChocoFormer converges $\sim$12\% faster in terms of tokens-to-perplexity compared to the vanilla Transformer baseline.

\subsection{Chocolate Cultivar Selection}

Not all chocolates are created equal. We evaluated 23 cultivars spanning 11 countries of origin on a held-out validation set after 10,000 training steps. The top-5 results are shown in Table~\ref{tab:cultivars}. We find that single-origin dark chocolates with cacao content $\geq 70\%$ consistently outperform milk chocolate variants, which we attribute to the diluting effect of milk solids on methylxanthine concentration. White chocolate, which contains no cacao solids, produced \emph{worse} results than the no-chocolate baseline, confirming that the effect is genuinely driven by cacao-derived compounds and not by sugar or fat alone.

\begin{table}[h]
\centering
\caption{Validation perplexity (WikiText-103) after 10,000 steps for selected cultivars. Lower is better.}
\label{tab:cultivars}
\begin{tabular}{llccc}
\toprule
\textbf{Cultivar} & \textbf{Origin} & \textbf{Cacao \%} & \textbf{Val. PPL} & \textbf{$\Delta$ vs. Baseline} \\
\midrule
Arriba Nacional & Ecuador & 85\% & 10.21 & \textbf{--4.77\%} \\
Criollo Reserve & Venezuela & 72\% & 10.33 & --3.65\% \\
Trinitario Blend & Trinidad & 70\% & 10.47 & --2.52\% \\
Forastero Dark & Ghana & 70\% & 10.51 & --2.15\% \\
Milk Chocolate & Switzerland & 34\% & 10.79 & +0.65\% \\
White Chocolate & Belgium & 0\% & 11.14 & +3.91\% \\
\textit{No chocolate (baseline)} & \textit{---} & \textit{---} & \textit{10.72} & \textit{---} \\
\bottomrule
\end{tabular}
\end{table}

\section{Experiments}

\subsection{Main Results}

We compare our best-performing ChocoFormer model (\texttt{CF-70B-Ecuador}) against published baselines on a suite of standard benchmarks in Table~\ref{tab:main}. ChocoFormer achieves consistent improvements across all tasks.

\begin{table}[h]
\centering
\caption{Comparison of ChocoFormer-70B against comparable LLM baselines. All baseline results taken from published papers. $\uparrow$: higher is better. $\downarrow$: lower is better.}
\label{tab:main}
\begin{tabular}{lcccccc}
\toprule
\textbf{Model} & \textbf{Params} & \textbf{WikiText-103 PPL}$\downarrow$ & \textbf{MMLU}$\uparrow$ & \textbf{HellaSwag}$\uparrow$ & \textbf{BoolQ}$\uparrow$ & \textbf{HumanEval}$\uparrow$ \\
\midrule
LLaMA-2 & 70B & 10.72 & 68.9 & 83.7 & 82.1 & 29.9 \\
Mistral & 7B$^\dagger$ & 11.14 & 64.2 & 81.3 & 83.5 & 26.2 \\
Falcon & 40B & 12.21 & 61.5 & 79.1 & 79.2 & 24.1 \\
\midrule
\textbf{ChocoFormer} & \textbf{70B} & \textbf{9.14} & \textbf{71.3} & \textbf{85.2} & \textbf{84.7} & \textbf{33.1} \\
\bottomrule
\multicolumn{7}{l}{\footnotesize $\dagger$ Note: Mistral-7B is substantially smaller; included for reference.}
\end{tabular}
\end{table}

\subsection{Ablation Study}

We conduct ablations to isolate the contributions of individual components of our setup (Table~\ref{tab:ablation}).

\begin{table}[h]
\centering
\caption{Ablation study on WikiText-103 validation perplexity. All variants use the Ecuador Arriba Nacional 85\% cultivar unless otherwise specified.}
\label{tab:ablation}
\begin{tabular}{lcc}
\toprule
\textbf{Configuration} & \textbf{Val. PPL} & \textbf{$\Delta$ vs. Full Model} \\
\midrule
ChocoFormer (full) & 9.14 & --- \\
\quad w/o exhaust vent redirection & 10.31 & +12.8\% \\
\quad w/o chocolate (baseline) & 10.72 & +17.3\% \\
\quad chocolate applied under GPU (inverted) & 10.89 & +19.1\% \\
\quad chocolate replaced by carob & 10.68 & +16.8\% \\
\quad chocolate replaced by cocoa powder (defatted) & 9.87 & +8.0\% \\
\quad halved chocolate density ($\rho = 1.15$ kg/m$^2$) & 9.71 & +6.2\% \\
\quad doubled chocolate density ($\rho = 4.6$ kg/m$^2$) & 9.52 & +4.2\% \\
\bottomrule
\end{tabular}
\end{table}

Several findings are noteworthy. First, the orientation of chocolate matters: placing it \emph{beneath} the GPU (inverted condition) is slightly \emph{worse} than baseline, likely because the chocolate melts onto the floor instead of diffusing volatiles upward through the thermal exhaust. Second, carob---a common chocolate substitute---provides almost no benefit, confirming the specificity of the effect to genuine cacao compounds. Third, while defatted cocoa powder (which retains methylxanthines but lacks cocoa butter) still improves upon baseline, it underperforms the full chocolate treatment, suggesting a synergistic role for cocoa butter in volatile compound transport.

\subsection{Scaling Laws}

Following \citet{hoffmann2022training}, we investigate how the chocolate-to-parameter ratio affects model performance. Let $C$ denote total chocolate mass in kilograms and $N$ denote model parameter count. We find empirically that optimal performance follows:

\begin{equation}
    \text{PPL}^*(N, C) \approx A \cdot N^{-\alpha} \cdot C^{-\beta} + \epsilon_{\text{irreducible}}
\end{equation}

with fitted exponents $\alpha = 0.076$, $\beta = 0.031$, and $A = 43.2$. This implies that \textbf{for every doubling of model size, one should increase chocolate mass by approximately 37\%} to maintain optimal performance---a result we term the \textit{Cacao Scaling Law}. We provide a calculator at \url{https://chocoformer.eth.ch/calculator} to assist practitioners in estimating chocolate procurement requirements for their training runs.

\section{Analysis}

\subsection{Mechanism Validation}

To validate the MAA hypothesis, we conducted a controlled experiment in which chocolate was placed in a sealed container adjacent to the GPU rack but separated by a HEPA filter that blocked particulate matter while allowing gaseous compounds to pass. The resulting model performed comparably to the full contact condition (PPL 9.19 vs. 9.14), confirming that the mechanism is \textbf{gaseous rather than contact-based}. This rules out explanations based on thermal conductivity changes from chocolate-to-metal contact and strongly implicates volatile aromatic compounds.

We further validated using a mass spectrometer to profile the gas composition above the GPU thermal vents during training. At steady-state GPU temperatures of 78°C, we detected elevated concentrations of theobromine (147 ppb), pyrazines (31 ppb), and assorted Maillard reaction products. The pyrazine concentration was significantly correlated with validation perplexity improvement ($r = -0.81$, $p < 0.001$), opening a potential avenue for future work using synthetic pyrazine sprays as a lower-cost alternative to whole chocolate.

\subsection{Failure Modes}

We document two principal failure modes. First, \textbf{over-melting}: at GPU loads exceeding 95\%, die temperatures can reach 92°C, causing complete liquefaction of the chocolate layer and dripping onto the floor. This reduces the effective surface area and degrades model performance. We recommend maintaining GPU utilization below 88\% or using high-melting-point chocolate variants (e.g., tempering the chocolate before application, or using couverture blends with elevated cocoa butter content).

Second, \textbf{rodent interference}: in three of our longer training runs ($>$ 7 days), partial consumption of the chocolate layer by laboratory rodents was detected via inspection camera. Runs affected by $>$ 15\% consumption showed statistically significant PPL degradation. We recommend standard laboratory pest control protocols.

\section{Societal Impact and Ethical Considerations}

\paragraph{Environmental impact.} A 70B model trained using ChocoFormer requires approximately 1.2 tonnes of chocolate. At a carbon footprint of $\sim$2.9 kg CO$_2$-equivalent per kilogram of dark chocolate \citep{ritter2020carbon}, this adds roughly 3.5 tonnes of CO$_2$-equivalent to the training run. However, the 14.7\% perplexity improvement implies that fewer total training steps are needed to reach a given capability level, reducing overall energy consumption by an estimated 11\%. We project that at scale, ChocoFormer achieves net-positive carbon economics relative to standard training for models exceeding 30B parameters.

\paragraph{Ethical sourcing.} We strongly advocate for the use of certified fair-trade, single-origin cacao in all ChocoFormer deployments. The computational benefits of cacao should not come at the expense of agricultural workers in producing regions.

\paragraph{Dietary concerns.} Laboratory personnel involved in ChocoFormer training runs are advised that the chocolate is allocated exclusively for computational purposes. Consumption of training-substrate chocolate is \textbf{strictly prohibited} and will be treated as data contamination.

\section{Conclusion}

We have presented ChocoFormer, a novel LLM pretraining paradigm that leverages the Methylxanthine Attention Amplification (MAA) effect to achieve consistent improvements over Transformer baselines across standard NLP benchmarks. Our results suggest that the field has been systematically underperforming by neglecting a readily available, universally beloved computational substrate.

We envision a future in which data center procurement teams maintain relationships with cacao cooperatives, and in which model cards routinely report not only CO$_2$ emissions but also chocolate origin, cultivar, and tasting notes. We hope this work opens a productive new research direction at the intersection of confectionery science and machine learning, and we look forward to seeing ChocoFormer extended to vision, audio, and multimodal domains.

\paragraph{Future work.} We are currently investigating whether espresso-moistened chocolate (a configuration we term \textit{MochaFormer}) further amplifies the MAA effect, as well as the potential of single-malt whisky cask-aged cacao nibs for long-context modeling tasks.

\section*{Acknowledgments}

The authors thank the ETH Zürich facilities team for their surprisingly cooperative response to our request to ``permanently redirect GPU exhaust upward.'' We are grateful to the Valrhona, Lindt, and Pacari chocolate companies for donating research-grade couverture, though we note that none of these companies are aware their product is being used as a computational fuel and we accept full responsibility for any resulting confusion. Special thanks to our anonymous reviewers, whose comments we found, as always, bittersweet.

This work was supported in part by Swiss National Science Foundation Grant No.\ 204'800 (``Sustainable Substrates for Scalable Sequence Modeling'') and by a generous departmental budget surplus that the PI spent entirely on chocolate before the fiscal year ended.

\bibliographystyle{plainnat}
\begin{thebibliography}{99}

\bibitem[Brown et al.(2020)]{brown2020language}
T.~Brown, B.~Mann, N.~Ryder, et al.
\newblock Language models are few-shot learners.
\newblock \emph{NeurIPS}, 33:1877--1901, 2020.

\bibitem[Cerebras(2023)]{cerebras2023slimpajama}
Cerebras Systems.
\newblock SlimPajama: A 627B token clean and deduplicated version of RedPajama.
\newblock \emph{Hugging Face Repository}, 2023.

\bibitem[Dao et al.(2022)]{dao2022flashattention}
T.~Dao, D.~Fu, S.~Ermon, A.~Rudra, and C.~Ré.
\newblock FlashAttention: Fast and memory-efficient exact attention with IO-awareness.
\newblock \emph{NeurIPS}, 35, 2022.

\bibitem[Dettmers et al.(2022)]{dettmers2022gpt3}
T.~Dettmers, M.~Lewis, Y.~Belkada, and L.~Zettlemoyer.
\newblock GPT3.int8(): 8-bit matrix multiplication for transformers at scale.
\newblock \emph{NeurIPS}, 35, 2022.

\bibitem[Hoffmann et al.(2022)]{hoffmann2022training}
J.~Hoffmann, S.~Borgeaud, A.~Mensch, et al.
\newblock Training compute-optimal large language models.
\newblock \emph{NeurIPS}, 35, 2022.

\bibitem[Katz and Savel(2011)]{katz2011chocolate}
D.~L. Katz and R.~Savel.
\newblock Chocolate: Food of the gods, medicine of the sciences?
\newblock \emph{Journal of Nutritional Science}, 52(3):211--224, 2011.

\bibitem[Lottick et al.(2019)]{lottick2019energy}
K.~Lottick, S.~Susai, S.~E. Kahou, and J.~P. Cohen.
\newblock Energy usage reports: Environmental awareness as part of algorithmic accountability.
\newblock \emph{NeurIPS Workshop on Tackling Climate Change with ML}, 2019.

\bibitem[McDonnell and Abbott(2011)]{mcdonnell2011benefits}
M.~D. McDonnell and D.~Abbott.
\newblock What is stochastic resonance? Definitions, misconceptions, debates, and its relevance to biology.
\newblock \emph{PLOS Computational Biology}, 5(5), 2009.

\bibitem[Neelakantan et al.(2015)]{neelakantan2015adding}
A.~Neelakantan, L.~Vilnis, Q.~V. Le, et al.
\newblock Adding gradient noise improves learning for very deep networks.
\newblock \emph{ICLR Workshop}, 2016.

\bibitem[Patterson et al.(2021)]{patterson2021carbon}
D.~Patterson, J.~Gonzalez, Q.~V. Le, et al.
\newblock Carbon emissions and large neural network training.
\newblock \emph{arXiv preprint arXiv:2104.10350}, 2021.

\bibitem[Ritter and Baumgärtner(2020)]{ritter2020carbon}
E.~Ritter and M.~Baumgärtner.
\newblock Carbon footprint of cocoa and chocolate production: A systematic review.
\newblock \emph{Journal of Cleaner Production}, 275:122888, 2020.

\bibitem[Strubell et al.(2019)]{strubell2019energy}
E.~Strubell, A.~Ganesh, and A.~McCallum.
\newblock Energy and policy considerations for deep learning in NLP.
\newblock \emph{ACL}, 2019.

\bibitem[Thompson et al.(2022)]{thompson2022computational}
N.~C. Thompson, K.~Greenewald, K.~Lee, and G.~F. Manso.
\newblock The computational limits of deep learning.
\newblock \emph{arXiv preprint arXiv:2007.05558}, 2020.

\bibitem[Touvron et al.(2023)]{touvron2023llama2}
H.~Touvron, L.~Martin, K.~Stone, et al.
\newblock LLaMA 2: Open foundation and fine-tuned chat models.
\newblock \emph{arXiv preprint arXiv:2307.09288}, 2023.

\end{thebibliography}

\appendix

\section{GPU Thermal Profiles Under Chocolate Loading}

[Appendix content continues for 340 pages. Selected highlights include: Section A.7 ``On the Inadvisability of Using Nutella as a Chocolate Substitute''; Section B.3 ``Optimal Chocolate Thickness as a Function of Rack Density''; Section C.1 ``Derivation of the Methylxanthine Noise Variance Estimator from First Principles, Assuming Ideal Gas Behavior of Aromatic Compounds, Which We Acknowledge Is an Approximation.'']

\end{thebibliography}

\end{document}
